Pathology foundation models turn whole-slide images into reusable representations from
which lightweight probes can recover grade, tissue phenotype, molecular labels and
clinical outcomes \cite{conch2023,virchow2023}. This capability has made candidate signals
much easier to produce than to interpret. A high pooled area under the receiver operating
characteristic curve (AUROC) or correlation can arise
because the representation captures the target, but it can also reflect grade composition,
institutional case mix, acquisition conditions, sampling choices or the definition of the
clinical endpoint. Consequently, the central problem is no longer only whether a frozen
embedding contains a statistically detectable association. It is whether the available
evidence is sufficient to qualify that association as transportable, to restrict it to a
specific context, or to leave it unsupported in the evaluated design.

This distinction matters because common validation practices collapse different scientific
questions into a single performance number. External transport asks whether a signal
retains direction in a different cohort without refitting. Confounder auditing asks whether
it contributes information beyond established clinical structure such as grade. Site
analysis asks whether a pooled result survives institutional heterogeneity. Representation
and sampling audits ask whether the interpretation depends on encoder, physical scale,
tile budget or sampling seed. Endpoint auditing asks whether an outcome association
survives a change in event definition rather than merely a change in name. Success on one
axis cannot substitute for missing evidence on another, and repeated estimates obtained
from the same patients and folds are sensitivity analyses rather than independent
replications.

We organized the study around three testable propositions. First, signals corresponding to
morphology directly expressed in routine histology should show stronger cross-cohort
transport than signals whose labels are molecular or outcome-derived. Second, an apparent
molecular or outcome association may narrow after grade, site, representation or endpoint
audits even when its pooled estimate is positive. Third, candidate signals should therefore
be assigned evidence states from the pattern of support across these dimensions, not ranked
or validated by one aggregate metric. These propositions define the logic evaluated here;
they do not imply that every target had an identical external cohort or that lack of support
in the frozen analysis establishes biological absence.

% Previous prose version retained for provenance:
% To test this logic, we assembled complementary evidence from five prostate-pathology data
% resources. NADT-Prostate provided source-cohort grade and phenotype evaluations
% \cite{nadt2021}; PANDA provided case images from Karolinska and Radboud for transfer without
% refitting \cite{panda2022}; PRECISE provided an additional session-level grade evaluation
% \cite{precise2026}; TCGA-PRAD provided molecular labels, site structure and explicitly
% distinguished recurrence endpoints \cite{tcgaprad2015,tcgacdr2018}; and the LEOPARD-derived
% analysis supplied a post-hoc recurrence signal for transfer \cite{leopard2026}. Across frozen
% CONCH and Virchow representations, we evaluated grade, tumor phenotype, PTEN loss, SPOP
% mutation, androgen-receptor activity and recurrence. The analyses linked cross-cohort
% transfer to held-out confounder increments, site-specific uncertainty, a correlated
% encoder--scale--tile--seed grid and endpoint-specific survival comparisons. This design
% allowed the same candidate to succeed on one evidentiary axis while remaining unresolved
% on another.

To test this logic, we assembled complementary evidence from five prostate-pathology data
resources. These were the prostate cohort treated with neoadjuvant androgen-deprivation
therapy (NADT-Prostate); the Prostate cANcer graDe Assessment (PANDA) challenge; the PRECISE
(PRostate Expert-annotated Contiguous IHC--H\&E Serial sEctions) dataset, where IHC denotes
immunohistochemistry and H\&E denotes hematoxylin and eosin; The Cancer Genome Atlas prostate
adenocarcinoma (TCGA-PRAD) cohort; and the
LEOPARD prostatectomy-slide challenge. Their roles are summarized below.

\begin{table}[!htbp]
\centering
\small
\setlength{\tabcolsep}{4pt}
\renewcommand{\arraystretch}{1.15}
\caption*{\textbf{Study data resources and their roles in signal qualification.} Each
resource supplied a distinct evidence component; no resource alone covered every
qualification dimension.}
\begin{tabular}{@{}p{0.21\textwidth}p{0.40\textwidth}p{0.30\textwidth}@{}}
\toprule
Resource & Evidence supplied & Role in the qualification design \\
\midrule
NADT-Prostate \cite{nadt2021} & Source-cohort images, grade and phenotype labels &
Source-cohort grade and phenotype evaluation \\
PANDA \cite{panda2022} & Karolinska and Radboud case images & Cross-cohort transfer
without refitting \\
PRECISE \cite{precise2026} & Session-level pathologist-assigned grade & Additional
grade evaluation \\
TCGA-PRAD \cite{tcgaprad2015,tcgacdr2018} & Molecular labels, source-site structure
and distinguished recurrence endpoints & Molecular, site and endpoint audits \\
LEOPARD \cite{leopard2026} & Post-hoc recurrence signal & Transfer to TCGA-PRAD \\
\bottomrule
\end{tabular}
\end{table}

Across frozen CONtrastive learning from Captions for Histopathology (CONCH) and Virchow
representations, we evaluated grade, tumor phenotype, phosphatase and tensin homolog (PTEN)
loss, speckle type BTB/POZ protein (SPOP) mutation, androgen-receptor (AR) activity and a
post-hoc recurrence risk signal. The analyses linked
cross-cohort transfer to held-out confounder increments, site-specific uncertainty, a
correlated encoder--scale--tile--seed grid and endpoint-specific survival comparisons. This
design allowed the same candidate to succeed on one evidentiary axis while remaining
unresolved on another.

The study makes three contributions. First, it operationalizes pathology foundation-model
signal qualification as a joint, target-level assessment of transport, confounding, site
behavior, representation stability and endpoint fidelity. The resulting categories---
transportable, context-sensitive and unsupported in the frozen design---describe the scope
of evidence rather than a league table of model performance. Second, it provides an
empirical evidence map across morphologic, molecular and outcome targets: grade and
phenotype showed the clearest transport; PTEN and androgen-receptor activity required
narrower interpretations; SPOP was unsupported in the fixed primary configuration; and
recurrence remained endpoint- and representation-conditioned. Third, it demonstrates why
pooled discrimination alone is insufficient: apparently positive probes can lose
incremental support, remain unresolved across sites or change interpretation with endpoint
and representation. The accompanying source-linked workflow makes those qualification
decisions auditable and reproducible. The contribution is therefore not another encoder or
a deployable classifier, but a concrete framework and empirical failure map for deciding
which foundation-model signals warrant further validation, under what conditions, and with
which claims.
